\chapter{2019年『数学分析A2』期中试题解答}
\vspace{4em}
\hspace{-2.16em}
\begin{minipage}{0.48\textwidth}
    \textbf{课程编号}：A0711172
\end{minipage}
\hfill
\begin{minipage}{0.48\textwidth}
    \textbf{考试时间}：$2019$年$4$月$29$日\quad $9:00\sim 11:00$\ AM
\end{minipage}
\begin{exercise}
    判断题【正确的记\faCheckCircleO，错误的记\faTimesCircleO】（每题2分，共10分）
\begin{enumerate}
    \item $\xi$为点集为点集$S$的聚点的充要条件：存在各项互不相同的数列$\ab{a_n}\subset S$，使得$\lim_{n\to\infty}{a_n}\xi$.\hfill【\qquad】
    \begin{solution}
    
    \end{solution}

    \item 若函数$f\ab(x)$在$\comm{a}{b}$上有界，则$f\ab(x)$在$\comm{a}{b}$上可积.\hfill【\qquad】
    \begin{solution}
    
    \end{solution}

    \item 若$\int_a^{+\infty}{f\ab(x)\d x}$收敛，且极限$\lim_{x\to+\infty}{f\ab(x)}$存在，则$\lim_{x\to+\infty}{f\ab(x)}=0$.\hfill【\qquad】
    \begin{solution}
    
    \end{solution}

    \item 若$\sum_{n=1}^{\infty}{a_n}$收敛，$\sum_{n=1}^{\infty}{b_n}$发散，则级数$\sum_{n=1}^{\infty}{\ab(a_n-b_n)}$发散.\hfill【\qquad】
    \begin{solution}
    
    \end{solution}

    \item 若级数$\sum_{n=1}^{\infty}{a_n}$收敛，$\sum_{n=1}^{\infty}{b_n}$发散，则级数$\sum_{n=1}^{\infty}{\ab(a_n-b_n)}$发散.\hfill【\qquad】
    \begin{solution}
    
    \end{solution}
\end{enumerate}
\end{exercise}

\begin{exercise}
    填空题（每题4分，共8分）
\begin{enumerate}
    \item （积分第二中值定理推论）若……23333，使得$\int_a^b{f\ab(x)g\ab(x)\d x}=\underline{\mbox{下划线就是这样用的}}$
    \item （阿贝尔？）……条件：
    \begin{enumerate}[(1)]
        \item 级数……
        \item 数列……
    \end{enumerate}
\end{enumerate}
\end{exercise}

\begin{exercise}
    计算下列各题（每题6分，共 30分，\textbf{所有题目要写出详细过程，只写结果不得分}）
    \begin{enumerate}
        \item 用定积分定义求极限$\lim_{n\to\infty}{n\ab[\frac{1}{\ab(n+1)^2}+\frac{1}{\ab(n+2)^2}+\cdots+\frac{1}{\ab(n+n)^2}]}$.
        \begin{solution}
            
        \end{solution}
    \end{enumerate}
\end{exercise}